% Ranked by how often this professor asks them. Each numbered line is worth
% roughly one mark, so the numbering doubles as a point checklist.
\subsection{How to write a proof for marks}\label{sub:proofhow}
\R{generic scheme}{
\begin{rcp}
\item Write the hypotheses and the goal (``Voraussetzung / Zu zeigen'').
\item Reduce w.l.o.g.\ (\emph{oBdA}) --- e.g.\ replace $f$ by $-f$ to fix a sign.
\item Build the auxiliary object: a set, a helper function, a subsequence. \emph{Most marks sit here.}
\item Justify that it exists (completeness / Bolzano--Weierstrass / Weierstrass max).
\item Run the $\eps$-argument or the contradiction.
\item Close with the statement and $\square$.
\end{rcp}
\hd{Cheap marks:} name every theorem you use and quote its hypotheses (``$f$ continuous on the \emph{compact} $[a,b]$, hence \dots''). Writing out the definitions of the given and the wanted object is often already half the marks.
}

\subsection{Tier 1 --- expect one of these}\label{sub:proofs1}
\F{Intermediate value theorem \de{Zwischenwertsatz} (9 P)}{
$f\in C([a,b])$, $c$ between $f(a)$ and $f(b)$ \ra\ $\exists x^*:f(x^*)=c$.
\begin{rcp}
\item oBdA $f(a)\le c\le f(b)$ (else consider $-f$).
\item $X:=\{x\in[a,b]\mid f(x)\le c\}$.
\item $X\ne\emptyset$ ($a\in X$) and $X\subseteq[a,b]$ is bounded \ra\ $x^*:=\sup X$ exists by completeness.
\item Choose $x_n\in X\cap[x^*-\frac1n,x^*]$ \ra\ $x_n\to x^*$ with $f(x_n)\le c$.
\item Sequential continuity \ra\ $f(x^*)=\lim f(x_n)\le c$.
\item Suppose $f(x^*)<c$. Then $x^*<b$ (as $c\le f(b)$). With $\eps:=c-f(x^*)>0$ continuity gives $\delta>0$ such that $|y-x^*|<\delta\Rightarrow|f(y)-f(x^*)|<\eps$.
\item Hence $f(y)<f(x^*)+\eps=c$ there, i.e.\ $(x^*,x^*+\delta)\cap[a,b]\subseteq X$.
\item That contradicts $x^*=\sup X$ \ra\ $f(x^*)=c$. $\square$
\end{rcp}
}
\F{Cauchy sequence converges \de{Vollst\"andigkeit von $\Rl$} (6 P)}{
\begin{rcp}
\item Cauchy \ra\ \textbf{bounded}: with $\eps=1$ get $N$, then $|a_n|\le\max\{|a_1|,\dots,|a_{N}|,|a_N|+1\}$.
\item Bolzano--Weierstrass \ra\ a convergent subsequence $a_{n_k}\to L$.
\item Cauchy: $\exists N$ with $|a_n-a_m|<\frac\eps2$ for all $n,m\ge N$.
\item Subsequence: $\exists N_1$ with $|a_{n_k}-L|<\frac\eps2$ for all $k\ge N_1$.
\item $N^*:=\max\{N,N_1\}$; since $n_k\ge k$, pick one index $n_k\ge N^*$.
\item Triangle inequality: $|a_n-L|\le|a_n-a_{n_k}|+|a_{n_k}-L|<\frac\eps2+\frac\eps2=\eps$ for all $n\ge N^*$. $\square$
\end{rcp}
}
\F{Mean value theorem \de{Mittelwertsatz} (3 P)}{
\begin{rcp}
\item Auxiliary function $g(x):=f(x)-\dfrac{f(b)-f(a)}{b-a}(x-a)$.
\item $g\in C[a,b]$, differentiable on $(a,b)$, and $g(a)=f(a)=g(b)$ --- the hypotheses of Rolle.
\item Rolle \ra\ $\exists\xi\in(a,b)$ with $0=g'(\xi)=f'(\xi)-\frac{f(b)-f(a)}{b-a}$. $\square$
\end{rcp}
}
\F{Rolle}{
\begin{rcp}
\item $f$ continuous on the compact $[a,b]$ \ra\ attains max and min (Weierstrass).
\item If both are attained at the endpoints, $f$ is constant \ra\ $f'\equiv0$, done.
\item Otherwise an extremum is interior, at some $\xi$ \ra\ Fermat \ra\ $f'(\xi)=0$. $\square$
\end{rcp}
}
\F{Bolzano--Weierstrass}{
Bisection. $(a_n)\subseteq[A,B]$. Halve the interval and keep a half containing infinitely many terms; repeat \ra\ nested intervals of length $\frac{B-A}{2^k}$. Pick indices $n_1<n_2<\dots$ with $a_{n_k}$ in the $k$-th interval. The nested-interval principle gives a single point $L$, and $|a_{n_k}-L|\le\frac{B-A}{2^k}\to0$. $\square$
}
\F{monotone $+$ bounded \ra\ convergent}{
$(a_n)$ increasing and bounded \ra\ $s:=\sup\{a_n\}$ exists (completeness). Given $\eps>0$, $s-\eps$ is not an upper bound \ra\ $\exists N:a_N>s-\eps$. Monotonicity \ra\ $s-\eps<a_N\le a_n\le s$ for all $n\ge N$ \ra\ $|a_n-s|<\eps$. $\square$
}

\subsection{Tier 2 --- likely}\label{sub:proofs2}
\F{extreme value theorem (Weierstrass)}{
\hd{(i) bounded:} otherwise $\exists x_n$ with $|f(x_n)|\ge n$; Bolzano--Weierstrass \ra\ $x_{n_k}\to x^*\in[a,b]$, continuity \ra\ $f(x_{n_k})\to f(x^*)$ finite --- contradiction.\\
\hd{(ii) attained:} $S:=\sup f$ exists by (i); pick $x_n$ with $f(x_n)\to S$; a convergent subsequence $x_{n_k}\to x^*$ gives $f(x^*)=S$ by continuity. $\square$
}
\F{continuous on a compact interval \ra\ uniformly continuous (Heine)}{
Suppose not: $\exists\eps_0>0$ and $x_n,y_n$ with $|x_n-y_n|<\frac1n$ but $|f(x_n)-f(y_n)|\ge\eps_0$. Bolzano--Weierstrass \ra\ $x_{n_k}\to x^*$, and then $y_{n_k}\to x^*$ too. Continuity \ra\ $|f(x_{n_k})-f(y_{n_k})|\to0$ --- contradiction. $\square$
}
\F{fundamental theorem of calculus}{
$F(x)=\int_a^xf$. Then $\frac{F(x+h)-F(x)}h=\frac1h\int_x^{x+h}f=f(\xi_h)$ for some $\xi_h$ between $x$ and $x+h$ (MVT for integrals). As $h\to0$, $\xi_h\to x$, so continuity gives $F'(x)=f(x)$. $\square$
}
\F{uniform limit of continuous functions is continuous}{
The $\frac\eps3$-argument:
$|f(x)-f(x_0)|\le|f(x)-f_N(x)|+|f_N(x)-f_N(x_0)|+|f_N(x_0)-f(x_0)|$.
Choose $N$ with $\|f_N-f\|_\infty<\frac\eps3$ (uniformity: the \emph{same} $N$ for all $x$ --- this is the whole point), then $\delta$ from continuity of $f_N$ at $x_0$. $\square$
}
\F{absolute convergence \ra\ convergence}{
Cauchy criterion for series: $\left|\sum_{k=n}^ma_k\right|\le\sum_{k=n}^m|a_k|<\eps$ for $m\ge n\ge N$. So the partial sums are Cauchy, hence convergent. $\square$
}
\F{Fermat}{
$x_0$ an interior local max: for $h>0$ small, $\frac{f(x_0+h)-f(x_0)}h\le0$; for $h<0$ it is $\ge0$. Both one-sided limits equal $f'(x_0)$, so $f'(x_0)=0$. $\square$
}
\F{differentiable \ra\ continuous}{
$f(x)-f(x_0)=\frac{f(x)-f(x_0)}{x-x_0}\cdot(x-x_0)\to f'(x_0)\cdot0=0$. $\square$
}

\subsection{Tier 3 --- one-liners you should still know}\label{sub:proofs3}
\F{convergent \ra\ bounded}{
$\eps=1$ gives $N$ with $|a_n-a|<1$ for $n\ge N$ \ra\ $|a_n|\le|a|+1$; the finitely many earlier terms are bounded by their maximum. $\square$
}
\F{convergent \ra\ Cauchy}{
$|a_n-a_m|\le|a_n-a|+|a-a_m|<\frac\eps2+\frac\eps2$ for $n,m\ge N$. $\square$
}
\F{the limit is unique}{
Assume $a_n\to a$ and $a_n\to b$ with $a\ne b$. Take $\eps=\frac{|a-b|}2$: then $|a-b|\le|a-a_n|+|a_n-b|<|a-b|$ --- contradiction. $\square$
}
\F{squeeze theorem}{
$b_n\le a_n\le c_n$ with $b_n,c_n\to L$: for large $n$, $L-\eps<b_n\le a_n\le c_n<L+\eps$. $\square$
}
\F{limit laws (sum, product)}{
Sum: $|(a_n+b_n)-(a+b)|\le|a_n-a|+|b_n-b|<\frac\eps2+\frac\eps2$.\\
Product: $|a_nb_n-ab|\le|a_n||b_n-b|+|b||a_n-a|$, and $(a_n)$ is bounded by some $C$ (convergent \ra\ bounded) \ra\ $\le C|b_n-b|+|b||a_n-a|<\eps$. $\square$
}
\F{product rule}{
$\frac{(fg)(x)-(fg)(x_0)}{x-x_0}=f(x)\frac{g(x)-g(x_0)}{x-x_0}+g(x_0)\frac{f(x)-f(x_0)}{x-x_0}$; let $x\to x_0$ and use $f(x)\to f(x_0)$. $\square$
}
\F{chain rule}{
Set $\varphi(y):=\frac{g(y)-g(y_0)}{y-y_0}$ for $y\ne y_0$ and $\varphi(y_0):=g'(y_0)$; then $\varphi$ is continuous at $y_0$ and $g(y)-g(y_0)=\varphi(y)(y-y_0)$ for all $y$. Insert $y=f(x)$, $y_0=f(x_0)$, divide by $x-x_0$, let $x\to x_0$. $\square$
}
\F{derivative of the inverse}{
$f^{-1}$ continuous and $f'(x_0)\ne0$: $\frac{f^{-1}(y)-f^{-1}(y_0)}{y-y_0}=\frac{x-x_0}{f(x)-f(x_0)}\to\frac1{f'(x_0)}$ as $y\to y_0$ (which forces $x\to x_0$). $\square$
}
\F{geometric series}{
$s_N=\sum_{k=0}^Nq^k$; then $s_N-qs_N=1-q^{N+1}$ \ra\ $s_N=\frac{1-q^{N+1}}{1-q}$. For $|q|<1$, $q^{N+1}\to0$ \ra\ $s_N\to\frac1{1-q}$. $\square$
}
\F{root test}{
$\limsup\sqrt[n]{|a_n|}=q<1$: pick $q<r<1$; for $n\ge N$, $\sqrt[n]{|a_n|}\le r$ \ra\ $|a_n|\le r^n$ \ra\ comparison with a convergent geometric series. If $q>1$ then $|a_n|\ge1$ infinitely often \ra\ $a_n\not\to0$. $\square$
}
\F{Cauchy--Hadamard}{
Apply the root test to $\sum a_n(x-x_0)^n$: $\limsup\sqrt[n]{|a_n||x-x_0|^n}=|x-x_0|\limsup\sqrt[n]{|a_n|}=\frac{|x-x_0|}R$, which is $<1$ exactly when $|x-x_0|<R$. $\square$
}
\F{comparison test}{
$0\le a_n\le b_n$ with $\sum b_n<\infty$: the partial sums of $\sum a_n$ are increasing and bounded above by $\sum b_n$ \ra\ monotone $+$ bounded \ra\ convergent. $\square$
}
\F{Leibniz criterion}{
$b_n\downarrow0$: $(s_{2n})$ increases, $(s_{2n+1})$ decreases, $s_{2n}\le s_{2n+1}$, and $s_{2n+1}-s_{2n}=b_{2n+1}\to0$ \ra\ nested intervals \ra\ both converge to a common $s$. $\square$
}
\F{monotone functions are Riemann integrable}{
$f$ increasing on $[a,b]$, equidistant partition of mesh $\frac{b-a}n$:
$U(f,P)-L(f,P)=\frac{b-a}n\sum_i\big(f(x_i)-f(x_{i-1})\big)=\frac{b-a}n\big(f(b)-f(a)\big)\to0$ by telescoping \ra\ integrable. $\square$
}
\F{uniform convergence allows swapping $\lim$ and $\int$}{
$\left|\int_a^bf_n-\int_a^bf\right|\le\int_a^b|f_n-f|\le(b-a)\|f_n-f\|_\infty\to0$. $\square$
}
\F{l'Hospital (idea)}{
With $f(x_0)=g(x_0)=0$, the Cauchy MVT gives $\frac{f(x)}{g(x)}=\frac{f(x)-f(x_0)}{g(x)-g(x_0)}=\frac{f'(\xi_x)}{g'(\xi_x)}$ for some $\xi_x$ between $x_0$ and $x$; $x\to x_0$ forces $\xi_x\to x_0$. $\square$
}
\F{Archimedean property}{
If $\N$ were bounded above, $s:=\sup\N$ would exist; $s-1$ is not an upper bound \ra\ $\exists n>s-1$ \ra\ $n+1>s$, contradicting that $s$ is an upper bound. Hence $\frac1n\to0$. $\square$
}
\F{characterisation of the supremum}{
$s=\sup M\iff$ (i) $x\le s\ \forall x\in M$ and (ii) $\forall\eps>0\ \exists x\in M:x>s-\eps$.\\
``\ra'': if (ii) failed for some $\eps$, then $s-\eps$ would be a smaller upper bound. ``$\Leftarrow$'': (i) makes $s$ an upper bound, (ii) makes it the least one. $\square$
}
\F{Cauchy--Schwarz}{
For every $t\in\Rl$: $0\le\|x+ty\|^2=\|x\|^2+2t\langle x,y\rangle+t^2\|y\|^2$. A real quadratic in $t$ that is never negative has discriminant $\le0$: $4\langle x,y\rangle^2-4\|x\|^2\|y\|^2\le0$. Equality forces a double root, i.e.\ $x+ty=0$ for some $t$. $\square$
}
\F{convex $\iff f''\ge0$ (differentiable case)}{
$f$ convex $\iff$ the chord slopes increase $\iff f'$ is monotonically increasing $\iff f''\ge0$ (monotonicity criterion). $\square$
}
\F{contraction: uniqueness of the fixed point}{
$|f(x)-f(y)|\le q|x-y|$ with $q<1$. Two fixed points give $|x^*-y^*|\le q|x^*-y^*|$ \ra\ $(1-q)|x^*-y^*|\le0$ \ra\ $x^*=y^*$. Existence: $x_{n+1}=f(x_n)$ satisfies $|x_{n+1}-x_n|\le q^n|x_1-x_0|$, hence is Cauchy. $\square$
}
\F{a convergent sequence has exactly one accumulation point}{
$a_n\to a$: every subsequence has indices $n_k\ge k$, so it also converges to $a$. Hence $a$ is the only subsequential limit. $\square$
}
